\section{背景}

\subsection{残差流与拒答方向}

残差流（residual stream）是 Transformer 各子层读写的共享隐藏表示。对同一层分别汇总有害提示与无害提示在回答起点处的残差，可以用两组均值之差构造行为方向。Arditi 等人把拒答与一维子空间联系起来，并通过删除或注入该方向改变模型是否拒答~\cite{arditi2024refusal}。该结果支持“拒答方向可干预”这一研究假设，但其 13 个模型不包含本文讨论的 Qwen3.8-27B，因而不能替代针对新架构的验证。

方向消融（directional ablation）从权重或激活中删除某个单位方向上的分量。若直接对输出投影矩阵做秩一修改，模型无需新的训练样本或梯度更新；但方向估计、作用组件和层权重仍决定实际效果。投影式消融进一步去除拒答方向在无害均值方向上的分量，目标是减少对无害表征的直接干扰。本文分析的实现以该直觉为起点，但其全量归一化（full normalization, FULL）和有限秩存储使最终更新不再是无损正交投影。

\subsection{低秩适配与行归一化}

低秩适配（Low-Rank Adaptation, LoRA）用两个小矩阵表示对基础权重的增量，从而避免为每个试验复制完整模型~\cite{hu2022lora}。秩一方向更新天然可写成列向量与行向量的乘积；该系统因此把每次消融写入 LoRA，并通过清零 B 矩阵恢复未消融模型。若需要保持完整更新后的行范数，逐行归一化会引入非线性，系统便先构造完整差值，再用有限秩分解近似它。LoRA 在这里是试验和导出的载体，不表示模型接受了常规参数高效微调。

\subsection{多目标黑盒优化}

树结构 Parzen 估计器（Tree-structured Parzen Estimator, TPE）根据历史试验提出新的参数候选。Optuna 提供 define-by-run 搜索空间和可扩展的试验管理~\cite{akiba2019optuna}；该系统使用其中的多变量 TPE，同时优化多个 scorer。对两个最小化目标，若一个试验无法在不恶化至少一个目标的情况下继续改进，它就属于 Pareto 前沿。这个集合保留了“更少关键词拒答”与“更小原模型分布漂移”之间的显式权衡。
